% Tabla C-T4. Generada automáticamente; no editar a mano.
% Requiere \usepackage{booktabs} (y graphicx si la tabla es ancha).
\begin{table}[htbp]
    \centering
    \caption{Pares adyacentes: correlación de Spearman y pendiente controlada entre separación orbital y desigualdad de tamaños en siete submuestras.}
    \label{tab:C-T4}
    \small
    \resizebox{\textwidth}{!}{%
    \begin{tabular}{lrrrrrrrrr}
    \toprule
    Submuestra & $N_{\rm pares}$ & $N_{\rm sis}$ & $\rho$ & $\rho$ inf. & $\rho$ sup. & Pendiente & IC 95\,\% inf. & IC 95\,\% sup. & $p$ \\
    \midrule
    Muestra amplia (P $<$ 100 d) & 1\,155 & 767 & 0.221 & 0.166 & 0.278 & 0.182 & 0.140 & 0.225 & $<0.001$ \\
    Ambos Rp $>$ 1 $R_\oplus$ & 1\,036 & 704 & 0.210 & 0.152 & 0.267 & 0.148 & 0.105 & 0.191 & $<0.001$ \\
    Ambos Rp $<$ 4 $R_\oplus$ & 980 & 658 & 0.187 & 0.123 & 0.244 & 0.142 & 0.103 & 0.181 & $<0.001$ \\
    Solo tránsito & 797 & 539 & 0.214 & 0.146 & 0.281 & 0.157 & 0.107 & 0.207 & $<0.001$ \\
    Sistemas de Kepler & 728 & 474 & 0.177 & 0.108 & 0.247 & 0.153 & 0.095 & 0.212 & $<0.001$ \\
    Sistemas con N = 2 & 505 & 505 & 0.217 & 0.128 & 0.302 & 0.156 & 0.109 & 0.203 & $<0.001$ \\
    Sistemas con N $\geq$ 3 & 650 & 262 & 0.202 & 0.125 & 0.275 & 0.255 & 0.176 & 0.334 & $<0.001$ \\
    \bottomrule
    \end{tabular}
    }
    \par\smallskip
    {\footnotesize\textit{Nota:} Pendiente de la separación orbital en el modelo U1, con errores agrupados por estrella. IC de rho por bootstrap (3000 remuestreos).}
\end{table}
